\documentclass[11pt,a4paper]{article}
\usepackage{fontspec}
\setmainfont{DejaVu Sans}[Scale=0.90]
\usepackage{amsmath,amssymb}
\usepackage[margin=2.2cm]{geometry}
\usepackage{booktabs,array,tabularx}
\usepackage{graphicx}
\graphicspath{{../figures/}}
\usepackage{enumitem}
\usepackage{titlesec}
\usepackage{parskip}
\titlelabel{\thetitle.\quad}
\titleformat{\section}{\normalfont\large\bfseries}{\thesection.}{0.6em}{}
\titleformat{\subsection}{\normalfont\normalsize\bfseries}{\thesubsection}{0.6em}{}
\titlespacing*{\section}{0pt}{15pt}{5pt}
\titlespacing*{\subsection}{0pt}{11pt}{4pt}
\setlist[itemize]{leftmargin=1.4em,itemsep=1pt,topsep=3pt}
% Preserve fonts/margins; permit a final line-breaking pass for long prose.
\setlength{\emergencystretch}{2em}
\begin{document}
\begin{center}
{\Large\bfseries Dégénérescences et fermeture cinétique}\\[4pt]
{\normalsize Phases, populations de bloc et mémoire}\\[2pt]
{\small Partie I --- 26 septembre 2026}
\end{center}

\section{Question et portée}
La note 20 constate que des interactions individuelles changent avec la
base choisie dans des sous-espaces numériquement dégénérés, tandis que
leurs sommes restent stables. Peut-on remplacer ces interactions par
leurs sommes et obtenir une dynamique autonome des populations de bloc ?

La réponse est négative en général. Les contre-exemples ci-dessous sont
finis et synthétiques. Ils ne démontrent pas une défaillance du modèle
de populations pour l'AlN réel. L'opérateur physique du matériau, ses
conventions de canaux et son régime de cohérence restent à établir.
Les critères employés relèvent de la théorie connue de réduction exacte
et de projection ; aucune nouveauté générale n'est revendiquée.

\section{Ce que les modules au carré ne déterminent pas}
Considérons deux espaces de dimension deux et une troisième branche
singleton, avec les amplitudes
\[
 T_+=\begin{pmatrix}1&1\\1&1\end{pmatrix},\qquad
 T_-=\begin{pmatrix}1&1\\1&-1\end{pmatrix}.
\]
Les matrices $|T_\pm|^2$ sont identiques, mais les spectres de
$T_\pm^\dagger T_\pm$ sont respectivement $\{0,4\}$ et $\{2,2\}$.
Une transformation unitaire indépendante sur les deux espaces préserve
ces spectres. Les deux tenseurs ne sont donc pas équivalents sous cette
action, malgré des modules au carré identiques.

Avec la rotation de Hadamard
$U=2^{-1/2}\left(\begin{smallmatrix}1&1\\1&-1\end{smallmatrix}\right)$,
\[
 |T_+U|^2=\begin{pmatrix}2&0\\2&0\end{pmatrix},\qquad
 |T_-U|^2=\begin{pmatrix}2&0\\0&2\end{pmatrix}.
\]
Les interférences nécessaires au changement de base sont absentes des
modules au carré. Ce contre-exemple établit une non-injectivité, pas
un théorème de généricité pour toute dimension.

Une matrice d'occupation se transforme avec ses cohérences. Ne conserver
que sa diagonale après chaque rotation perd de l'information.
En revanche, dans une base enregistrée et sous une approximation
incohérente déclarée, les modules au carré peuvent suffire à définir
des taux si les incidences, facteurs statistiques et comptages sont connus.

\section{Critère exact de fermeture des populations de bloc}
Soit un système linéaire fini
\[
 \dot x=-Lx,\qquad N=Bx,
\]
où $B$ est de rang ligne plein. Une équation $\dot N=-GN$ valable pour
\emph{tout} état initial existe si et seulement si
\[
                         BL=GB.
\]
La nécessité résulte de la dérivée à l'instant initial ; la suffisance
s'obtient par substitution. Cette condition concerne un générateur
spécifié, non la seule symétrie des fréquences.

Dans le modèle réversible de populations, posons
\[
 W_{ii}=n_i^0(1+n_i^0)>0,\quad
 M=\sum_e\lambda_e\nu_e\nu_e^T,\quad L=MW^{-1},\quad \lambda_e\ge0.
\]
Avec
\[
 \Sigma=BWB^T,\qquad H=WB^T\Sigma^{-1},\qquad P=HB,
\]
on a $BH=I$ et l'unique candidat réduit est
\[
 G=BLH=BMB^T\Sigma^{-1}.
\]
Il s'agit d'une compression. La fermeture exige en plus
\[
                         BL(I-P)=0.
\]
Pour un bloc dégénéré de dimension $d_A$ et de variance par mode $w_A$,
$\Sigma_{AA}=d_Aw_A$ et $(HN)_i=N_A/d_A$.
La variance d'une population totale n'est pas obtenue en remplaçant
$n_i^0$ par cette population totale dans la formule de Bose.

\section{Contre-exemple conservatif et positif}
Prenons quatre modes d'énergies réduites
\[
 (\log6,\log6,\log2,\log3),\qquad
 W=\operatorname{diag}(6/25,6/25,2,3/4),
\]
et deux événements distincts $p_1\leftrightarrow a+b$,
$p_2\leftrightarrow a+b$. Ils sont exactement résonants.
Le facteur de Bose commun vaut $3/5$. Comparons les préfacteurs
$\Gamma=(2,2)$ et $\Gamma=(3,1)$ dans les unités de temps de l'exemple.
La somme des poids et la matrice comprimée sont identiques :
\[
 G=\begin{pmatrix}
 5&-6/5&-16/5\\
 -5&6/5&16/5\\
 -5&6/5&16/5
 \end{pmatrix}.
\]
Le premier modèle ferme, le second non. Pour $h=(1,-1,0,0)^T$,
\[
 Bh=0,\qquad -BLh=(-5,5,5)^T
\]
dans le second modèle. Deux états ayant les mêmes populations de bloc
peuvent donc avoir des dérivées différentes.
Une perturbation suffisamment petite conserve la positivité des
occupations. Les deux modèles conservent l'énergie.

Même une préparation uniforme dans le bloc parent ne suffit pas :
elle développe ensuite un déséquilibre caché. Treize vérifications
symboliques et numériques testent ce cas ; pour l'état sélectionné,
l'erreur de compression sur les populations atteint $4.43\,\%$ à
$t=0.5$ dans les unités de cet exemple. Ce nombre n'est pas une erreur
de mesure ni une prédiction sur l'AlN.

\section{Ce que l'élimination exacte ajoute}
En coordonnées d'entropie, soit $C=C^\dagger\ge0$ et la décomposition
en variables retenues $a$ et cachées $h$ :
\[
 C=\begin{pmatrix}A&K\\K^\dagger&D\end{pmatrix}.
\]
L'élimination exacte donne
\[
 \dot a(t)=-Aa(t)+\int_0^t
 K e^{-D(t-s)}K^\dagger a(s)\,ds-Ke^{-Dt}h(0).
\]
La mémoire a un signe positif dans cette convention ; le terme
d'état initial caché a un signe négatif.
Pour $\Re z>0$, l'opérateur de Schur est
\[
 zI+A-K(zI+D)^{-1}K^\dagger.
\]
Si $h(0)=0$, la compression reproduit la pente initiale, mais
$\ddot a(0)=(A^2+KK^\dagger)a(0)$. La fermeture pour tous les états
requiert $K=0$.

La positivité implique $K\ker D=0$. Elle n'élimine pas les invariants
du système complet et n'assure pas l'existence d'un inverse statique.
Onze contrôles supplémentaires, dont un calcul à 80 chiffres,
vérifient la mémoire et l'action du modèle de la note 19.
Ils montrent aussi qu'un pôle conservé dominant peut masquer une
erreur de $25\,\%$ sur le canal relaxant dans une norme relative globale.

\section{Retour au triplet réel de l'AlN}
Pour le triplet enregistré $(1,13,28)$ de la note 20, les constantes
de force permettent de reconstruire les amplitudes complexes avant
la prise du module au carré. Une contraction indépendante des indices
normaux retrouve l'export natif avec une erreur globalement normalisée
inférieure à $8\times10^{-16}$.
La transformation complète entre bases donne $2.15\times10^{-15}$ ;
en ne permettant que des rotations dans les blocs numériques,
l'erreur relative de Frobenius vaut $2.46\times10^{-14}$.

La revue indépendante retrouve ces valeurs. Les erreurs relatives par
canal peuvent néanmoins atteindre $2.28\times10^{-11}$.
La contraction partage le tenseur réciproque, les constantes et la
normalisation amont : elle ne vérifie pas toute la physique du matériau.

La revue a aussi détecté un défaut général du regroupement fréquentiel :
une chaîne de fréquences proches pouvait créer des groupes recouvrants
et une transformation non unitaire. Le triplet étudié n'est pas affecté.
Le programme portable rejette désormais les chaînes ambiguës, vérifie
la partition et l'unitarité. Huit régressions passent ; les résultats du
triplet restent inchangés. L'ancien algorithme et son contre-exemple
sont conservés.

Un seuil fréquentiel ne certifie pas une dégénérescence physique exacte.
Les phases reconstruites sont celles du tenseur à trois indices entrants,
pas encore celles de canaux de décroissance normalisés et comptés.

\section{Conséquences pour la suite}
Il faut distinguer quatre produits : événements de populations dans
une base enregistrée ; compression invariante ; dynamique fermée
certifiée ; action sur des matrices d'occupation.
La covariance d'un changement de coordonnées n'est pas une symétrie
physique du générateur fixé. Un tenseur complexe ne définit pas, à lui
seul, une cinétique irréversible.

La prochaine étape est de déclarer un canal et une approximation
cinétique, puis de vérifier séparément orientation, conjugaison,
comptage, conservation et fermeture. Aucun opérateur complet de l'AlN,
aucune convergence matérielle et aucune validation expérimentale ne
sont établis ici.

\section{Sources et traçabilité}
Le dossier \texttt{theory/aln/degenerate\_action/} conserve les quatre
branches indépendantes, la seconde génération, les quatre revues
adversariales, les scripts, les résultats et les échecs.
Les identités finies ont été auditées ; aucune vérification formelle
par assistant de preuve n'a été effectuée.

Le cadre bibliographique comprend Ovchinnikov et al.,
\emph{CLUE}, arXiv:2004.11961 ; Simoncelli, Marzari et Mauri,
Phys.\ Rev.\ X \textbf{12}, 041011 (2022) ; Trushechkin,
Phys.\ Rev.\ A \textbf{103}, 062226 (2021) ; et la théorie de
projection et de mémoire, avec sources et limites d'accès précisées
dans le registre bibliographique.
\end{document}
